\chapter{怎样证明看不见的东西存在}

\begin{kaichang}
夏天傍晚，妈妈在客厅里喷了一点花露水，然后去厨房忙了。你坐在离她三米远的沙发上，没过半分钟，一股清凉的香味就钻进了你的鼻子。

停下来想一想这件事有多奇怪：花露水明明还留在桌上那瓶子里，谁也没有端着它凑到你面前，你却闻到了它。\textbf{有什么东西从瓶子里跑出来，穿过了三米的空气}——可你睁大眼睛，什么也看不见。

它存在吗？当然存在，你的鼻子就是证人。可鼻子不会说话，它拿不出照片，也没法把"香味"放到秤上称一称。这就引出了整门科学最头疼、也最迷人的问题：\textbf{一样东西看不见、摸不着、拍不到，我们凭什么说它存在？}这一章，我们就来学这套"给看不见的东西定罪"的手艺。学完之后你会发现，从花露水的香味，到比细菌还小几万倍的原子，用的其实是同一套逻辑。
\end{kaichang}

\section{看见的和看不见的}

先把话说清楚：科学研究的对象，绝大多数是肉眼看不见的。分子看不见，细菌看不见，风看不见，温度看不见，就连"氧气"这种东西，你把眼睛瞪得再圆也看不见它。

那科学是不是就只能靠猜？恰恰相反。科学有一套极其实用的办法：\textbf{不直接看它，改看它留下的痕迹}。

这办法你早就会用了。冬天早上起床，你看见窗户玻璃上有一层水汽，立刻知道"昨晚外面很冷"——你没站在窗外感受过，但水汽就是证据。放学回家闻到楼道里有饭香，你知道谁家在做晚饭。看见树叶在动，你知道"起风了"——风这东西，从古到今没有任何人亲眼看见过它，我们看见的全是风推着动的东西：树叶、旗子、水面上的波纹。

科学家把这种思路分成了两层：

\begin{itemize}[itemsep=2pt]
\item \textbf{直接观察}：现象本身能被你的感官（必要时借助简单仪器放大感官）接收。比如天平的指针偏了、试纸变色了、温度计液柱上升了。
\item \textbf{间接证据}：研究对象本身看不见，但它的\emph{效应}可以被测量。比如没人看见过电子，但电子打在荧光屏上会让屏发光——光看得见，于是电子"被看见了"。
\end{itemize}

关键在于：一条间接证据可能靠不住，也许有别的解释。玻璃上的水汽也可能是屋里太潮；饭香也可能是别人家的抽油烟机串了味。所以科学真正的功夫，在于\textbf{把许多条彼此独立的间接证据串成链条}，让所有证据指向同一个解释，并且排得掉其他解释。证据链越完整，"看不见的东西存在"这句话就越硬气——硬气到比"我亲眼所见"还可靠，因为眼睛会骗人，而彼此独立的测量一起骗人的机会要小得多。

\section{侦探的六件工具}

要给看不见的东西定罪，得有称手的工具。下面这六样，是物质侦探工具箱里最常用的，你以后会在每一章反复用到它们。

\textbf{第一件：质量。}天平是最老实的证人，它不参与任何争论。一样东西来了还是走了，让质量说话。把一块铁钉放在潮湿空气里，过几天再称，它变重了——多出来的质量从哪来？一定有看不见的东西（空气里的成分）加入了铁钉。第 1 章说过，"有没有新物质生成"是判断化学变化的线索，而质量的增减往往是这条线索最硬的版本。

\textbf{第二件：体积。}气体看不见，但它占据空间。给气球充气，气球鼓起来；堵住注射器的口往里推活塞，越推越费劲——里面那团看不见的东西在抵抗。体积测量让"空气占据了空间、还能被压缩"这件事变成数字，而不再是一句空话。

\textbf{第三件：密度。}质量除以体积，得到密度。它像物质的"身材密码"：同样大小的一块，铁比铝沉得多，铝比木头沉得多。两块外表一模一样的金属块，一称一量算密度，立刻知道里面装的料不一样。考古学家用密度鉴别金币有没有掺假，用的就是这把尺子。

\textbf{第四件：温度。}温度告诉你粒子"闹腾"的平均程度（第 4 章会把这句话讲透）。化学反应常常伴随着温度的变化：生石灰遇水烫手，冰袋里的硝酸铵溶于水变凉。温度计就是伸向粒子世界的探针。

\textbf{第五件：酸碱性。}用 pH 试纸或 pH 计测量。pH 是个数字，大致表示溶液里"氢离子"这种粒子的浓度高低——数字小（小于 7）偏酸，数字大（大于 7）偏碱，7 附近算中性（这个标准只在常温附近成立，第 32 章会细讲它的边界）。柠檬汁、肥皂水、胃液、雨水，各有自己的 pH，一测便知底细。

\textbf{第六件：光谱。}这是最神奇的一件，值得多说两句。让一束光通过棱镜或光栅，会展开成彩虹；而如果把某种元素烧亮、通电点亮，它发出的光展开后不是连续的彩虹，而是几条特定位置、特定颜色的亮线——每种元素的亮线排列都不一样，就像指纹。氖灯的红色、钠灯的黄光，都是"指纹"里的一条线。1868 年，天文学家在太阳的光谱里发现了一组对不上任何已知元素的亮线，于是宣布：太阳上有一种地球上还没见过的元素，并用希腊语"太阳"给它起名"氦"。二十七年后，人们才在地球上找到氦气——\textbf{一种元素，先在九千多万公里外被光谱"看见"，才在地球上被找到}。至于光谱线为什么能成为指纹，答案藏在原子内部电子的排布里，那是第 9 章的故事。

注意这六件工具的共同套路：\textbf{没有一件是"直接看见"研究对象的，全是在测量它产生的效应。}看不见没关系，测得到就行。

\section{让看不见的东西露出马脚}

工具在手，来看一场科学史上最漂亮的"间接定罪"。

\begin{zhengju}
1827 年，英国植物学家布朗把花粉撒在水面上，用显微镜观察，发现花粉颗粒在永不停歇地乱动：忽左忽右，走走停停，方向毫无规律。他起初以为是花粉"活着"，可换成磨细的石头粉末、玻璃粉末——任何足够小的颗粒，只要泡在水里，全都这样乱动。不是生命，那是什么在推它们？

这个谜题悬了将近八十年。期间有人猜是水温不均造成的水流，有人猜是光照，可实验排除了这些解释：恒温、避光，颗粒照样乱动，而且颗粒越小、水温越高，动得越欢。

1905 年，爱因斯坦给出了一个大胆的定量解释：水是由无数看不见、不停乱撞的水分子组成的，花粉颗粒悬在水里，每时每刻都被四面八方的水分子撞击。颗粒足够大时，各方向的撞击互相抵消，看不出动静；可花粉只有几微米，某一瞬间左边撞来的分子恰好比右边多几个，颗粒就被"顶"了一下——下一瞬间又换另一个方向。那些看似无规则的乱动，其实是\textbf{数不清的水分子撞击的统计涨落}。爱因斯坦还算出了颗粒平均应该走多远、跟水温、颗粒大小、水的黏稠程度的具体数学关系。

理论漂亮，可它预言的"水分子"谁也没见过，凭什么信？1908 年起，法国物理学家佩兰做了一连串苦功夫实验：他设法做出大小均匀的树脂颗粒，在显微镜下一颗一颗地追踪几百个颗粒的位置，记录下成千上万个数据点。结果发现：颗粒的实际运动跟爱因斯坦公式的预言严丝合缝。更妙的是，佩兰还从另一个完全不同方向的测量——数出悬浮颗粒随高度的分布——推算出了"每摩尔物质包含多少个粒子"这个数，即阿伏伽德罗常数，得到的值约为 $6\times 10^{23}$，与其他方法的结果一致。

到这时，连最顽固的怀疑者都低头了。反对"原子分子真实存在"的化学家奥斯特瓦尔德公开承认：布朗运动和佩兰的实验，"使最谨慎的科学家也有权谈论物质由不连续粒子构成"。注意这条证据链的结构：\textbf{看得见的花粉乱动 → 只能用水分子撞击解释（排除了水流、温度不均等替代解释）→ 数学预言被精确测量验证 → 多条独立路径算出同一个数}。没有任何人"看见"过水分子，可这条链比任何肉眼都可靠。佩兰因此获得 1926 年诺贝尔物理学奖。
\end{zhengju}

同样的套路，还两次改写了原子模型的历史。一次是阴极射线：抽成真空的玻璃管两端加上高压电，管中会出现一束看不见的"射线"，它能让荧光屏发光、能被磁铁掰弯、还能推动小风轮转——汤姆孙据此断定，这束射线是一股带负电的微小粒子流，也就是电子。另一次是金箔实验：卢瑟福让一束氦核粒子射向极薄的金箔，本以为会全部笔直穿过，结果极少数粒子被弹得倒飞回来，好像炮弹打在纸上被反弹——他由此推断，原子的质量和正电荷集中在一个小得惊人的核里。这两个实验的精彩细节，第 6 章会专门讲；这里你只要记住共同点：\textbf{射线和原子核都看不见，但荧光、偏转、反弹这些看得见的效应，把它们的存在钉死了}。

\section{好实验的规矩}

单个观察可以是巧合，一次测量可以是失误。要让"间接证据"变成"证据链"，实验必须守几条规矩。这些规矩不是科学家的洁癖，而是几百年被骗出来的教训。

\textbf{规矩一：能重复。}一个实验只有你做过一次、别人怎么都重复不出来，那它什么也证明不了。1989 年，两位化学家宣布在常温下的电解池里实现了"冷核聚变"，轰动全球；可全世界几十个实验室照着做，谁也重复不出来，结论随即被撤回。重复是科学给所有结论设的第一道关卡。

\textbf{规矩二：有对照。}想知道"肥料让水稻增产"，光在田里撒肥料、秋天丰收了，说明不了什么——也许今年雨水本来就好。正确的做法是把田分成两块，一切条件相同，一块施肥一块不施，这块不施肥的田就叫\textbf{对照组}。两组之间唯一不同的那个条件，叫\textbf{变量}。好实验的理想状态是：\emph{一次只改一个变量}，否则结果出来你根本不知道是谁的功劳。

\begin{shiyan}
\textbf{称一称"跑掉的空气"：蒸发对照实验}

材料：两个完全相同的小杯子（或小碗）、自来水、一小块保鲜膜、橡皮筋、一台厨房电子秤（没有秤也可以用记号笔）。

步骤：
\begin{enumerate}[itemsep=1pt]
\item 两个杯子各倒入同样多的水（半杯即可）。
\item 一杯敞口，另一杯用保鲜膜蒙住杯口、用橡皮筋扎紧。
\item 分别称重并记录（用记号笔的话，在杯壁上标出初始水位）。
\item 把两杯并排放在窗台或室内同一位置，静置两到三天。
\item 再次分别称重（或观察水位）。
\end{enumerate}

观察点：敞口杯质量减少了多少？封口的杯子呢？如果封口杯质量几乎不变而敞口杯明显变轻，少掉的质量去了哪里？

安全提示：只用清水，常温放置，没有任何危险。注意杯子别放在会被碰倒的地方。

这个实验里，封口杯就是对照组，"敞口还是封口"是唯一的变量。敞口杯变轻这条证据，逼你承认：有一部分水变成了看不见的水蒸气，穿过空气跑掉了——你看不见它离开，但秤看见了。这正是本章的核心手艺。
\end{shiyan}

\textbf{规矩三：该盲就盲。}人是很容易被自己的期待欺骗的动物。医生如果知道哪个病人吃的是真药、哪个吃的是长得一模一样的淀粉片（安慰剂），他的观察和判断难免偏心；病人知道自己吃了"新药"，也常会感觉好些。所以现代药物试验采用\textbf{双盲}：病人和直接观察的医生都不知道谁吃了什么，分组名单由第三方保管，直到统计结束才揭开。你可能觉得"我没那么容易受骗"——但证据恰恰表明，人人都有期待效应，盲法是用来保护诚实人的。

\textbf{规矩四：相关不等于因果。}这是最容易翻车的一条。统计发现：冰淇淋销量高的月份，溺水事故也多。是冰淇淋导致溺水吗？显然不是——两件事背后的共同原因是"天热"：天热吃冰淇淋的人多，下河游泳的人也多。冰淇淋销量和溺水数只是\emph{相关}，它们之间没有因果箭头。再比如"公鸡打鸣"和"太阳升起"天天相伴，可你绝不会认为太阳是被鸡叫出来的。判断因果，需要对照实验排除其他解释，或者找到环环相扣的机制链条——光有两条一起起伏的曲线远远不够。这条规矩请记牢，你今后在新闻、广告和保健品宣传里会天天用到它。

\section{测量永远带着"大约"}

还有一件必须坦白的事：\textbf{任何测量都不是绝对准确的}。尺子有最小刻度，秤有分度值，你的眼睛估读最后一位时也在猜。说"这杯水是 $200\,\mathrm{mL}$"，老实的说法是"大约是 $200\,\mathrm{mL}$，误差在几毫升以内"。这个"误差范围"，科学上叫\textbf{不确定度}。

写测量结果时，位数本身就是信息。"$2\,\mathrm{g}$""$2.0\,\mathrm{g}$""$2.00\,\mathrm{g}$"是三句不同的供述：第一句表示你只敢保证到个位，第二句表示你测到了 0.1 克这一档，第三句表示你的秤能分辨 0.01 克。这些"有把握的数字位数"叫\textbf{有效数字}。多写一位是吹牛，少写一位是浪费——秤明明能读到 0.01 克，你只记"2 克"，等于把好仪器用成了糙仪器。

\begin{qiaoliang}
\textbf{测不出小的，就测一堆再除}——这是对付"太小"的祖传智慧，用数字感受一下：

一粒米的质量大约是 $0.02\,\mathrm{g}$。家用厨房秤的分度值通常是 $1\,\mathrm{g}$。把一粒米放上去，读数是 $0$——不是米没有质量，而是它比秤的"分辨力"小了五十倍，秤根本看不见它。

换个办法：数出 500 粒米，一起称，读数约 $10\,\mathrm{g}$。这台秤的误差大约是 $\pm 1\,\mathrm{g}$，对 $10\,\mathrm{g}$ 来说相对误差约 $10\,\%$。于是每粒米的质量 $=10\,\mathrm{g}\div 500=0.02\,\mathrm{g}$，误差也摊薄到约 $10\,\%$。\textbf{单个测不准，就测五百个再平均}。

这个朴素的思想，正是科学家对付原子的办法：一个原子轻到 $10^{-23}$ 克量级，任何秤都无能为力，但只要数出天文数字那么多个一起称，再除以个数，一个原子的质量就出来了。问题在于——怎么"数"出天文数字个原子？这正是第 5 章的主角"摩尔"要解决的难题，而佩兰数布朗运动颗粒的实验，就是最早把这件事做成的尝试之一。
\end{qiaoliang}

\section{把测量写成符号}

测量有了数字，还得有一套大家通用的写法，否则你的"一拃"和我的"一拃"对不上账。这套写法就是科学界的符号语言，规则不多，但全世界都遵守：

\begin{itemize}[itemsep=2pt]
\item \textbf{数字带单位}。"铁钉重 $5$"是句不完整的话——$5$ 什么？克？千克？写全了：$m = 5.2\,\mathrm{g}$。单位不是装饰，单位错了，答案就错了，桥梁和火箭都为此出过事故。
\item \textbf{很大和很小的数用科学记数法}。一滴水里的分子数约为 $1.7\times 10^{21}$ 个；一个氢原子的质量约为 $1.7\times 10^{-24}\,\mathrm{g}$。指数 $21$ 表示小数点后面跟着 21 位，指数 $-24$ 表示小数点前面隔着 24 位。这套符号让天文数字和微观数字都能装进一行，第 5 章你会天天和它打交道。
\item \textbf{认真的时候，把不确定度也写出来}。比如 $m = (2.00 \pm 0.05)\,\mathrm{g}$，意思是：真值大概率落在 $1.95$ 到 $2.05$ 克之间。这个"$\pm$"不是谦虚，是测量结果的一部分——缺了它，别人无法判断你的证据有多硬。
\end{itemize}

你可能注意到了：到这一步，"花露水的香味"这件小事，已经被翻译成一串带单位、带误差范围的符号记录——"敞口杯质量从 $152.3\,\mathrm{g}$ 变为 $149.1\,\mathrm{g}$，封口杯从 $151.8\,\mathrm{g}$ 变为 $151.7\,\mathrm{g}$"。看不见的东西，就这样在纸上留下了供词。

\section{证据链也有边界}

这套"看效应不看本体"的手艺极其强大，但它有几条明确的边界，越过边界就会翻车。

\textbf{第一，仪器测到的永远要经过"解释"才能成为证据。}温度计液柱上升，我们解释为"变热了"；可液柱上升本身只是酒精在膨胀，"膨胀代表温度"这个对应关系是人为校准出来的。仪器越精密，中间隔的解释层越多。读仪器之前，先问一句：它实际测的到底是什么？

\textbf{第二，测量都有量程，超出量程的结论无效。}分度值 $1\,\mathrm{g}$ 的秤称一粒米，读数 $0$ 不代表质量是 $0$；普通显微镜看不见病毒，不代表病毒不存在——它只是小过了光学显微镜的极限（光波的波长本身决定了极限），直到电子显微镜出现才被"看见"。"没测到"和"不存在"是两回事，科学史上一半的进步来自有人造出了更灵敏的"秤"。

\textbf{第三，温度、pH 这类量，测的是"集体表现"而不是"某个粒子的状态"。}一杯 $25\,^{\circ}\mathrm{C}$ 的水，里面的水分子有的快有的慢，温度只是亿万个分子运动的平均结果。拿温度去描述单个分子，就像拿"全市平均工资"去描述某个具体的人——概念上用错了地方。

\textbf{第四，间接证据再强，也保留"被更好的解释推翻"的可能。}这不是软弱，而是科学的诚实。布朗运动把"分子存在"钉得非常死，但科学永远给新证据留着门——只不过证据链越完整，翻案需要的证据就得越硬。

\begin{wuqu}
\textbf{误解："眼见为实，看不见的东西就不存在。"}——照这个标准，空气不存在，细菌不存在，细菌和病毒被发现之前人类得传染病的历史就成了一笔糊涂账。真实情况恰恰相反：科学最重大的发现，几乎全是"看不见的东西"——原子、电子、基因、电磁波。判断存在与否的标准从来不是"肉眼能否看见"，而是"能否产生可重复、可测量、能排除其他解释的效应"。反过来也要警惕另一个极端："仪器屏幕上的数字就是终极真相。"仪器有量程、有误差、有解释层，试纸变色可能受到溶液本身颜色的干扰，再贵的仪器也只能在它的设计范围内说话。侦探既不相信"没看见就是没发生"，也不相信"仪器说的都对"——他只相信\textbf{彼此独立、相互印证、经得起重复}的证据链。
\end{wuqu}

\section{回到那瓶花露水}

现在，重新审理开场那桩"香味失踪案"。

花露水里有酒精和有香味的物质，它们的分子在液体表面不停地运动，总有一些撞开邻居的拉扯，跑进空气——这就是蒸发，第 4 章会详细拆解这个过程。跑出来的分子在空气里乱撞乱走（和布朗推着我们看到的花粉运动是同一套机制），几分钟内就散布到整个房间。几个幸运的香味分子飘进你的鼻腔，撞上嗅觉细胞，引发一连串化学信号传进大脑——你"闻到了"。

整个过程中，你的眼睛没有提供任何证据。但证据链完整无缺：敞口杯实验里的质量损失，说明有东西离开液体进了空气；温度越高香味散得越快，符合"温度越高分子越闹腾"的规律；换个房间关上门，香味就过不来，说明它依赖空气的流动和分子的扩散。每一条都可测量、可重复、排除了"心理作用"之外的所有替代解释——\textbf{看不见，但铁证如山}。

\section{这门手艺的用武之地}

你可能觉得"间接证据""对照组"这些规矩离生活很远，其实它们每天都在决定你的日子。你吃的每一种药，上市前都要闯过随机、双盲、有安慰剂对照的临床试验——正是本章这几条规矩，把"感觉有效"和"确实有效"分开，无数听起来动人的疗法倒在了这一关，这是好事。你手机上的空气质量指数，来自仪器对看不见颗粒物的持续测量；气象台的气温纪录、医院里的化验单、考古队的碳十四测年，全都是"测效应、推本体"的手艺。甚至广告里"九成使用者都说好"这类话术，你只要问一句"对照组呢"，就能当场拆穿大半。

而这一章的真正用意，是为整本书铺路。接下来我们要讲的几乎每一件东西——原子、电子、化学键、细胞里的分子机器——都看不见。你之所以能理直气壮地学它们，就是因为手里已经有了这套定罪手艺：\textbf{看效应，设对照，算误差，串证据链}。下一章我们先用它解决一个近在眼前的问题：你身边的物质，到底有多少是"纯"的？然后第 5 章攻克"数原子"的难题，到了第 6 章，你会看到"原子存在"这个今天看来天经地义的结论，当年是怎样被定比定律、被阴极射线、被金箔上反弹回来的粒子，一步一步\textbf{逼}出来的——那是这套手艺最辉煌的战例。

\begin{tiaozhan}
\textbf{为什么"多测几次取平均"会更准？}（跳过不影响主线）

假设每次测量的误差是随机的：有时偏大，有时偏小，平均而言不偏不倚。把 $N$ 次独立测量平均起来，随机误差会怎样变化？统计学的回答是：平均值的典型误差大约按 $1/\sqrt{N}$ 缩小。测 1 次误差是 $\pm 1$，测 4 次取平均误差降到约 $\pm 1/2$，测 100 次降到约 $\pm 1/10$。想再精确一位，测量次数就得乘以 100——这就是为什么精密实验又贵又慢。佩兰追踪几百个颗粒、记录成千上万个位置，本质上就是在用庞大的 $N$ 换取压倒性的精度。这个 $1/\sqrt{N}$ 规律有个漂亮的推论：大数目的随机事件，相对涨落反而小。一摩尔气体有 $6\times 10^{23}$ 个分子在乱撞，涨落只占 $\sqrt{N}$，即 $10^{12}$ 分之一量级——小到完全测不出，所以我们摸到的气压、水温稳定得像铁一样。\textbf{宏观世界的稳定，正是微观世界疯狂乱动的统计结果。}
\end{tiaozhan}

\section*{练一练}

\begin{sikaoti}
\item 冬天湖面结冰，冰面下的水却没有结冰。你能想到哪些"看不见的证人"帮助鱼类活过冬天？至少说出两条可以测量的效应（提示：想想温度和密度）。
\item 某广告说："本市常喝某某凉茶的人群中，感冒发生率比全市平均值低 $30\,\%$。"列出至少两种"凉茶有效"之外的解释，并设计一个能排除这些解释的实验（写出对照组、变量和要不要用盲法）。
\item 画一幅粒子示意图：敞口烧杯里的水和它上方的空气，用小点表示水分子，画出"有分子跑出液面、在空气中乱撞"的图景；再画同一烧杯加上密封盖之后的情形。图旁注明：点和轨迹是人为表示法，不是分子的真实模样。
\item 用分度值 $1\,\mathrm{g}$ 的厨房秤，估测一枚回形针的质量。设计测量方案并算出大致的相对误差（一枚回形针约 $0.5\,\mathrm{g}$）。如果换成分度值 $0.01\,\mathrm{g}$ 的电子天平，方案可以怎样简化？
\item 布朗运动中，为什么越小的颗粒乱动越明显？用"四面八方分子撞击的统计涨落"解释（提示：想想挑战层里的 $\sqrt{N}$：撞击次数少时，各方向抵消得干净还是不干净？）。
\item 这一章讲了六件"侦探工具"。针对"证明空气里有氧气参与铁生锈"，你会选哪几件工具、分别测什么？把你的证据链画成流程图：每个方框写测量结果，箭头写"由此推断"。
\end{sikaoti}
